% Transcription par lecture visuelle directe, sans OCR.
% Statut : transcrit ; relecture indépendante à effectuer.
% Source : PDF original, vue 118, côté gauche, page manuscrite 236, après le chapitre neuvième.
\chapter{Première suite de la préfecture de M. Franchet}
\begin{center}
Arrestation et captivité de MM. Franchet et du Coin.
\end{center}

Lorsque Napoléon croyait apercevoir dans les Congrégations comme un boulevard élevé pour défendre les droits du St Siège contre les envahissements dont il les menaçait, il ne se trompait pas entièrement. À la vérité la Congrégation de Lyon ne prit en cette qualité aucune part active dans la lutte ouverte
% Source : PDF original, vue 118, côté droit, page manuscrite 237.
entre la puissance temporelle et la puissance spirituelle. La plupart des congréganistes se bornaient comme les simples fidèles à gémir dans le secret de leurs cœurs et à supplier le Seigneur de donner la paix à son Église ; mais ce fut cependant parmi eux que se trouvèrent les hommes qui prirent en mains, dans notre ville, la cause du chef de l'Église persécuté.

Un petit nombre de congréganistes s'entendit et se concerta pour soutenir, en exposant au besoin leur liberté et même leur vie, la catholicité attaquée dans sa base. Si par prudence ils agissaient en dehors de la Congrégation, c'était néanmoins d'elle qu'ils recevaient l'impulsion de ce dévouement caractère propre des véritables congréganistes ; et la Congrégation dans laquelle ils avaient puisé l'esprit de foi qui les animait ne les aurait certainement pas désavoués.

J. C. qui a enfanté son Église sur le calvaire et qui montre constamment en sa faveur une sollicitude toute maternelle a soin de lui ménager pour les circonstances difficiles des âmes généreuses et privilégiées. Il dispose de loin ces âmes à marcher à sa suite dans la voie des tribulations et à affronter sans crainte toute espèce de danger pour la défense de la religion. Il les couvre du bouclier de la foi ; les remplit du courage des martyrs et de la présence d'esprit des confesseurs. C'est ainsi que dans tous les siècles cette sainte cause a toujours trouvé des défenseurs dignes d'elle. Grâces en soient rendues à notre bon maître, Messieurs ; aujourd'hui c'est dans la Congrégation de Lyon qu'il veut bien choisir ceux qu'il destine à cette noble mission ; il a désigné lui-même parmi ses membres les athlètes qu'il appelle à combattre pour sa
% Source : PDF original, vue 119, côté gauche, page manuscrite 238.
gloire ; il a enflammé leurs cœurs de son amour ; il leur a inspiré pour son service une énergie toute nouvelle ; il veut que, comme autrefois les Machabées ils raniment par l'ardeur de leur zèle le courage et la foi de leurs frères.

Au nombre de ces âmes d'élite la Congrégation pouvait compter son préfet M. Franchet ; ensuite cet autre congréganiste dont déjà nous avons parlé, notre vertueux ami de sainte et courageuse mémoire Bertaud du Coin.

\section*{Établissement d'une correspondance avec le pape}
% Mention marginale de la page manuscrite 238 ; placée au début du récit de la mission.
Il était indispensable pour pouvoir servir efficacement la cause de la religion de parvenir à se mettre en rapport avec le souverain pontife. Il fallait pouvoir déposer à ses pieds l'expression du dévouement des catholiques et braver toutes les rigueurs de l'inquisition de la police et des persécutions qui pouvaient en être la suite pour faire connaître à la France les ordres et les décisions qui émaneraient du Saint Siège. Mais où trouver quelqu'un d'assez dévoué pour aller au milieu des dangers tracer la ligne que devrait parcourir la correspondance qu'il faut établir à cette fin ? La Providence avait destiné du Coin à cette périlleuse mission ; il s'offrit généreusement pour la remplir et il s'en acquitta avec autant de dextérité que de courage. Deux fois il a fait à pied le voyage de Savonne ; il est parvenu jusqu'aux pieds du St Père et a eu le bonheur de verser quelques consolations dans son âme affligée ; il a pu l'entretenir de cette multitude de prêtres et de simples fidèles qui invariablement unis à l'Église catholique se tenaient plus que jamais attachés dans la bonne comme dans la mauvaise fortune à cette pierre angulaire posée par J. C. lui-même. Chère Congrégation de Lyon ! Si dans ces entrevues
% Source : PDF original, vue 119, côté droit, page manuscrite 239.
si précieuses mais environnées de tant de dangers, du chef de l'Église avec un de ses membres, il n'a pas pu être nominativement question de toi, oh ! n'en doute pas ! tu étais très certainement dans la pensée intime de celui des tiens qui remplissait cette glorieuse mission, car il te connaissait et savait que tu étais alors comme tu l'es encore aujourd'hui, dévouée de cœur et d'âme au St Siège apostolique, et que tous ceux qui t'appartiennent auraient été disposés à se lever à l'envi pour faire de leurs corps un rempart au pontife persécuté. (Si nous n'apercevons sur la brèche qu'un petit nombre d'entre eux, c'est qu'il n'y a pas eu place pour tous. Tous les cœurs des congréganistes étaient animés des mêmes sentiments. Ils se seraient tous bien volontiers pressés autour des combattants pour vaincre ou succomber avec eux.)

Du Coin établit sur toute la ligne de Savonne à Paris des points intermédiaires de correspondance. Elle passait par Chambéry et par Lyon, d'où elle était dirigée sur Paris. Le dernier anneau de cette chaîne était M. d'Astros, vicaire général de Paris. Par lui tous les évêques de France pouvaient recevoir toutes les communications qui émanaient du chef de l'Église. Ce fut par ce moyen que furent successivement transmises toutes les pièces officielles de la correspondance de Sa Sainteté avec Napoléon, et notamment la bulle d'excommunication qui fut fulminée par le St Père contre le persécuteur et tous ceux qui lui servaient d'instrument.

C'était sans doute déjà avoir beaucoup obtenu que d'être parvenu à faire connaître tous ces documents aux pasteurs
% Source : PDF original, vue 120, côté gauche, page manuscrite 240.
des peuples, mais ce n'était point encore assez ; il était à désirer qu'ils parvinssent aussi à la connaissance des simples fidèles ; il importait d'éclairer le public, de faire paraître la vérité au grand jour, afin d'imposer silence et de donner un éclatant démenti au mensonge, car le gouvernement qui tenait la presse muselée sous son joug de fer pouvait sans être contredit par personne donner à la résistance du pape les couleurs les plus odieuses. Il eût donc été bien utile de faire imprimer et publier toute cette correspondance ; il faut en convenir, en présence des cent yeux toujours ouverts de la police impériale, cette entreprise pouvait paraître hardie et presque audacieuse ; mais l'intérêt de la religion l'exigeait ; nos amis n'hésitèrent pas ; tout fut combiné et calculé avec autant de prudence que de zèle ; un imprimeur dévoué qui avait en sa possession une presse clandestine, se chargea de tout préparer pour pouvoir livrer sous peu une édition de 1200 exemplaires.

\section*{Impression de la correspondance du pape}
% Mention marginale de la page manuscrite 240 ; placée au début du récit des envois.
Cet imprimeur était au moment de commencer son travail lorsque survint la catastrophe qui devait enlever à la Congrégation et priver de leur liberté les principaux membres. Les paquets de Savonne parvenaient à Chambéry à M. l'abbé Rey, celui-ci les adressait à Lyon à M. du Coin qui à son tour les transmettait à M. l'abbé Perrault à Paris. Pour les y faire arriver avec plus de sûreté M. du Coin les remettait à Lyon à M. Franchet ; M. Franchet les expédiait à M. Varey, congréganiste de Lyon fixé à Paris, et M. Varey devait les porter à M. Perrault qui à son tour les mettait entre les mains de M. d'Astros. La police avait
% Source : PDF original, vue 120, côté droit, page manuscrite 241.
vainement cherché jusqu'ici à pénétrer le mystère dont cette ligne de correspondance était environnée ; les mesures avaient été si bien prises par M. l'abbé Perrault, l'âme et le principal moteur de cette sainte ligue qu'il avait rendu jusqu'alors toutes ces investigations inutiles. Mais lorsque la bulle d'excommunication parvint à la connaissance du despote, elle redoubla sa fureur, car si on prétend aujourd'hui que grâce aux lumières du siècle on ne craint plus les foudres de Rome, il n'en est pas moins vrai qu'elles paraissent encore redoutables à ceux contre lesquels elles sont lancées. La rage qu'éprouva l'empereur lorsqu'il se vit frappé par cette sentence du vicaire de J. C. prouverait au besoin la terreur qu'elle lui inspirait.

L'événement a prouvé qu'il avait bien raison de redouter ces traits, quoiqu'ils fussent lancés sur lui par la main débile d'un vieillard. Dieu permit alors que ces armes toutes spirituelles ne fussent pas sans force même en ce qui concernait les affaires de ce monde. Jusqu'à présent Napoléon monté sur le char de la victoire traîné rapidement par ses aigles triomphantes, s'était promené en vainqueur dans presque toutes les capitales de l'Europe. À peine la voix menaçante de Rome s'est-elle fait entendre que déjà le char de la victoire semble se ralentir ; si toutefois il le conduit encore à Moscou, ce sera pour s'y briser et précipiter le héros au fond d'un abîme. Encore quelques jours et les frimats de la Russie, l'abdication de Fontainebleau, les désastres de Waterloo et l'île Ste Hélène viendront démontrer aux yeux des moins
% Source : PDF original, vue 121, côté gauche, page manuscrite 242.
clairvoyants, quels effets le ciel permet quelquefois qu'une sentence d'excommunication produise.

Jusques à quand les peuples et les rois auront-ils donc des yeux pour ne pas voir ! Oh ! du moins, Messieurs, nous qui sommes congréganistes, sachons admirer ici les effets de la Providence de notre Dieu ; adorons sa justice, bénissons sa miséricorde ; que notre confiance se ranime, et que notre foi soit semblable à celle d'Abraham, car comme lui nous avons vécu dans un temps bien fertile en miracles.

Pardonnez-moi ces réflexions que la considération de ces grands événements fait échapper de ma plume presque malgré moi, je retourne à Napoléon. Dans sa colère, il veut absolument qu'on lui découvre ceux par l'intermédiaire desquels cette bulle d'excommunication est parvenue afin d'en tirer vengeance. Il lance tous les limiers de sa police comme une meute de chiens acharnés pour dépister les fidèles amis du St Père qui ont été assez hardis pour en mettre sous les yeux et presque signifier à celui qui fait trembler l'Europe les arrêts de la cour de Rome.

\section*{Découverte de la correspondance}
% Mention marginale de la page manuscrite 242 ; placée au début des arrestations de Paris.
La position officielle de M. d'Astros lui avait fait un devoir de se déclarer ouvertement ; il est arrêté. On parvient à savoir que c'est par M. l'abbé Perrault que lui ont été transmis les actes émanés du souverain pontife. M. Perrault est donc aussi jeté dans une prison. Pour achever de découvrir entièrement les secrets dont elle commençait à tenir quelques fils, la police place des agents dans la loge du portier de M. Perrault, avec ordre d'arrêter tous ceux qui se présenteraient pour le demander. M. Varey était dans l'ignorance la plus
% Source : PDF original, vue 121, côté droit, page manuscrite 243.
complète de tout ce qui se passait, il venait de recevoir une lettre de M. du Coin, pour M. Perrault et étant fort empressé de la remettre, il donne en plein dans le piège ; au moment où il se présente chez le portier il est arrêté et la police s'empare en même temps de la dépêche qui contenait de nouveaux documents de Savonne et un billet de M. Franchet qui lui avait servi d'enveloppe. Voilà par cet incident la police parfaitement au courant de la manière dont les pièces parvenaient de Lyon à Paris. Aussitôt le télégraphe est mis en mouvement pour donner l'ordre d'arrêter de suite MM. Franchet et du Coin. Ceci se passait le 10 janvier 1811, et le 11 au matin cet ordre fut exécuté. La police envahit simultanément le domicile de nos deux amis et mit en même temps la main et sur leurs personnes et sur tous leurs papiers. Ces papiers furent envoyés à Paris, au ministère de la police générale. Sans qu'on se donnât la peine de les examiner à Lyon, il ne se trouva rien de relatif à la Congrégation chez M. du Coin ; mais chez M. Franchet, on s'empara ainsi que nous l'avons appris seulement depuis sa délivrance, d'une copie du règlement de la Congrégation, celle-là même qui dans le temps avait été destinée à être remise au commissaire général de police. On n'y fit du reste qu'une bien médiocre attention ; cette pièce passa comme inaperçue et s'est retrouvée plus tard dans les papiers de M. Franchet, sans avoir jamais donné lieu à aucune question ; mais une pièce dont la saisie pouvait avoir les conséquences les plus graves échappa par bonheur ou plutôt par une disposition toute bienfaisante de la Providence divine, aux argus de l'inquisition impériale.

% Source : PDF original, vue 122, côté gauche, page manuscrite 244.
\section*{Arrestation de MM. Franchet et du Coin}
% Mention marginale des pages manuscrites 243–244 ; placée au début du récit du document caché.
C'était le tableau de la dernière organisation de la Congrégation. On y lisait en tête le nom de M. Franchet en sa qualité de préfet ; venait ensuite celui de M. du Coin comme premier assistant ; suivaient ceux des autres dignitaires et des congréganistes occupant les différents emplois, puis ceux des membres de la section du zèle, enfin les noms de tous les congréganistes divisés en trois sections avec leurs chefs en tête. Ce tableau se trouvait dans le secrétaire de M. Franchet dans un tiroir placé sous la case du milieu dont il était séparé par une coulisse ; c'est ce qu'on appelle la cave du secrétaire. Cette cave n'était masquée en aucune manière ; elle se trouve assez ordinairement dans tous les secrétaires de ce genre. Il paraît inconcevable que la pensée ne soit pas venue aux hommes de la police d'examiner si cette case n'existait pas au secrétaire qu'ils visitaient. Pour expliquer un semblable défaut d'attention de la part d'hommes de cette espèce, je crois que nous n'avons qu'un seul moyen, c'est d'admettre ici l'intervention de notre auguste patronne. Elle leur a fasciné les yeux, ces yeux de lynx si accoutumés à tout voir et à tout pénétrer. Notre divin maître avait permis qu'ils fussent ainsi aveuglés, car s'il avait destiné la Congrégation à être battue par l'orage, s'il voulait aussi faire passer au creuset de la persécution pour les éprouver comme l'or ceux d'entre les congréganistes qui lui étaient les plus agréables, il voulait en même temps que cette société subsistât afin qu'elle se préparât à accomplir plus tard les desseins qu'il avait sur elle. Or la prise de cette fatale liste eût été suivant toute apparence son arrêt de mort ; les deux congréganistes compromis se trouvant être les principaux chefs de la Congrégation, on n'aurait pas manqué d'en conclure que la société tout entière était complice ;
% Source : PDF original, vue 122, côté droit, page manuscrite 245.
on aurait mis la main sur les plus influents d'entre nous, on aurait employé les menaces et les persécutions pour disperser les autres, tout rapport et toute union particulière entre les congréganistes devenaient à l'avenir à peu près impossibles. Admirons donc, Messieurs, cette bonté, cette attention toute maternelle de notre bon Sauveur qui daigne ainsi veiller lui-même sur l'existence de notre Congrégation. Ne le voyons-nous pas ici ordonner à ses anges de nous garder dans toutes nos voies de peur que notre pied ne vienne à heurter contre quelque pierre ? Comment nos cœurs ne se dilateraient-ils pas en présence de ce Sauveur pour savourer avec délices les sentiments de la reconnaissance qui doit les pénétrer ?

Dans le même moment où on mettait MM. Franchet et du Coin en état d'arrestation, on se transportait également au domicile de M. Alexis de Noailles pour l'arrêter aussi. M. de Noailles habitait alors Lyon depuis plusieurs mois ; on le soupçonnait (et c'était avec raison) d'être au nombre des complices de nos deux amis, j'ignore comment le gouvernement avait été informé qu'il avait pris part à ce pieux complot. S'il se fût trouvé à Lyon au moment fatal, nul doute qu'il eût partagé la captivité des deux victimes de la tyrannie. Mais le Seigneur avait choisi parmi nous les héros qui devaient faire rejaillir sur la Congrégation l'honneur et le mérite de leurs souffrances. M. de Noailles ne devait pas être du nombre. Sans se douter le moins du monde du danger qui le menaçait il était parti pour Genève il y avait peu de jours, pour y faire une visite à un parent. Genève appartenait alors à la France. L'ordre de l'arrêter y parvint bientôt, mais prévenu à temps
% Source : PDF original, vue 123, côté gauche, page manuscrite 246.
par les soins d'un ami, il put se mettre en sûreté en passant en Suisse.

Nos deux prisonniers ignoraient absolument tous les deux si la mesure qui les avait frappés ne concernait qu'eux seuls, ou s'ils avaient des compagnons d'infortune. Ils furent mis en prison séparément et au secret le plus rigoureux. Dans la nuit qui suivit leur arrestation tous deux furent jetés dans des chaises de poste en la compagnie de gendarmes, et entraînés à Paris, l'un par la route du Bourbonnais, l'autre par celle de la Bourgogne. M. Franchet fut enchaîné pour ce voyage ; l'intervention de M. de Satonnay, maire de Lyon, son parent, préserva M. du Coin de cette rigueur. Il eut trouvé pendant le voyage l'occasion de se dérober à la rage du tyran ; une servante d'auberge voulait le faire évader par une porte de derrière, mais M. du Coin avait donné sa parole ; plus encore par principe que par honneur il en était esclave. Au risque de tout ce qui pourrait en advenir, il refusa donc cette chance de salut.

En arrivant à Paris, M. Franchet fut conduit à la prison de Ste Pélagie, et M. du Coin à celle de la Force. Là pendant cinq mois ils furent tenus tous les deux à un secret tellement rigoureux que leurs familles elles-mêmes furent pendant longtemps privées de pouvoir apprendre au moins ce qu'ils étaient devenus.

La police n'avait que trop pénétré nos secrets ; toutefois elle n'avait pas encore tout découvert et malgré toute son astuce elle ne put guère aller plus loin ; elle soupçonna cependant M. l'abbé Rey sans pouvoir en acquérir la preuve d'être pour quelque chose dans l'envoi des pièces officielles. Ce simple soupçon suffit pour motiver l'arrestation de ce vertueux ecclésiastique, mais on se borna à lui donner pour prison le séminaire de Chambéry,
% Source : PDF original, vue 123, côté droit, page manuscrite 247.
un factionnaire placé à la porte de sa chambre était chargé de le garder à vue.

Nous allons pour la première fois trouver ici l'occasion de parler de M. Recorbet de ce prêtre d'une vertu si angélique qui devait quelques années plus tard avoir avec la Congrégation des rapports si intimes, il était alors supérieur du petit séminaire de l'Argentières. Quoiqu'il fût étranger à la transmission des pièces, il fut véhémentement soupçonné d'être l'un des acteurs de cette sainte conspiration. Il faut bien convenir qu'il le méritait un peu par l'éclat de ses vertus et son inébranlable attachement au Saint Siège. S'il n'était pas coupable d'action, il l'était certainement au moins d'intention.

\section*{Arrestation de M. Recorbet}
% Mention marginale de la page manuscrite 247.
Les gendarmes de Trévoux reçurent l'ordre de l'arrêter et de le conduire à Lyon. Ils se transportèrent au petit séminaire pour exécuter leur mandat ; ils y arrivèrent le soir et signifièrent leurs ordres. M. Recorbet reçut avec cette douceur qu'il possédait si bien, les combla de politesses et les engagea à se reposer tranquillement la nuit et à attendre le lendemain de grand matin pour l'emmener. On voyait sur sa physionomie briller l'expression d'une douce joie ; il la puisait dans le souvenir de ces paroles prononcées par la vérité éternelle : Vous serez bienheureux lorsque à cause de moi les hommes vous chargeront d'opprobres, qu'ils vous persécuteront et qu'ils diront contre la vérité toute sorte de mal de vous. Réjouissez-vous alors et faites éclater votre joie, parce qu'une grande récompense vous attend dans le ciel.

Les séminaristes éprouvaient la douleur la plus vive en voyant ainsi arracher de cette sainte maison un supérieur qu'ils chérissaient
% Source : PDF original, vue 124, côté gauche, page manuscrite 248.
tous comme un père ; ils avaient de la peine à contenir les sentiments d'exaspération dont ils étaient animés ; ils auraient voulu s'insurger pour le délivrer des mains de ceux qui venaient se saisir de sa personne, mais le St prêtre s'opposa à tous ces projets ; à l'exemple du Sauveur du monde qui calma l'ardeur de ses apôtres et guérit l'oreille de Malchus, il contint aussi l'élan d'un zèle inconsidéré qui aurait pu attirer les plus grands malheurs sur cette maison et même sur le diocèse. Ce digne supérieur marchait à la suite de Jésus ; comme lui il se livrait à ses ennemis, et redisait après lui : remettez votre épée dans le fourreau, n'opposez aucune résistance. Est-ce que ce calice qu'on me présente, n'est-ce pas Dieu mon père qui me l'a préparé. Puis-je refuser de le boire ?

M. Recorbet fut conduit à Nancy. Après quelques jours on lui rendit un peu de liberté ; le ministre se borna à le retenir dans la ville qu'on lui donna pour prison.

À Lyon plusieurs autres personnes furent l'objet des recherches de la police ; on les manda à sa barre et leur fit subir des interrogatoires qui durent paraître fort extraordinaires à plusieurs d'entre elles. Quelque habile que pût être ordinairement cette police impériale, dans cette circonstance elle donna complètement à gauche. Pendant qu'on interrogeait et qu'on vexait ainsi des hommes inoffensifs parce qu'ils avaient eu quelques rapports de société avec M. de Noailles, qu'on leur parlait de pièces dont ils ignoraient jusqu'à l'existence, et du pape sur le compte duquel ils étaient parfaitement ignorants, les véritables complices ne furent pas inquiétés le moins du monde.

\section*{Publication des pièces officielles du Saint-Père}
% Mention marginale de la page manuscrite 248 ; placée au début de la poursuite de l'impression.
Comme on aperçoit ici avec évidence l'action de cette volonté suprême qui suspend l'effet des passions des hommes, comme
% Source : PDF original, vue 124, côté droit, page manuscrite 249.
la violence des flots de la mer, aussitôt que les uns et les autres ont atteint la limite qu'elle-même leur a tracée ! Je crois entendre encore la voix du Sauveur du monde qui au moment où il se livre pour notre salut entre les mains des pécheurs défend avec autorité d'attenter à la liberté de ses apôtres : Puisque vous tenez celui que vous cherchez, laissez aller ceux-ci. Cet ordre émané d'une bouche divine est exécuté à la lettre ; de même ici Dieu avait marqué dans ses conseils éternels ceux d'entre nos amis qui devaient le glorifier par leurs souffrances, il les a livrés à leurs persécuteurs ; mais aussitôt après l'esprit d'aveuglement vient obscurcir les yeux des agents du despote, afin que ceux qui avaient été associés aux confesseurs dans la pensée de défendre la religion restassent placés de manière à pouvoir mettre la dernière main à l'œuvre commencée. Ainsi que nous l'avons déjà dit on allait imprimer la correspondance du pape au moment où nos amis furent arrêtés ; bien loin qu'on considérât les rigueurs de la persécution comme un motif pour suspendre une publication qui pouvait en augmenter encore la sévérité, on crut au contraire que le retentissement que les mesures de violence donnaient à cette affaire était un motif pour imprimer cette correspondance à un plus grand nombre d'exemplaires. Au lieu de douze cents, il fut décidé qu'il y en aurait deux mille. En dépit de la police, l'impression eut lieu sans être troublée en aucune manière. La distribution de l'ouvrage dans toutes les mains catholiques la suivit de près, et tous les efforts du gouvernement impérial ne purent pas empêcher ce recueil de se répandre dans toute la France. Jamais on n'a pu découvrir ni l'imprimeur ni les distributeurs. Il régnait
% Source : PDF original, vue 125, côté gauche, page manuscrite 250.
un tel esprit d'union et de sage discrétion parmi les catholiques que quoique cet ouvrage fût à Lyon entre les mains de tous les fidèles, et que nous ayons vu nous-mêmes des personnes assez imprudentes pour en faire la lecture en parcourant les rues, la police, malgré la rage véritablement infernale avec laquelle elle le cherchait, n'a jamais pu parvenir à s'en procurer un seul exemplaire.

\section*{Détails sur la captivité de nos amis}
% Mention marginale de la page manuscrite 250.
Pendant que ces choses se passaient à Lyon, nos pauvres captifs supportaient avec une noble fermeté tout le poids de la fureur de Napoléon. Ils furent l'un et l'autre interrogés tour à tour à diverses reprises par les sicaires de la police, et firent paraître dans leurs réponses cette simplicité et cette grandeur d'âme qui furent toujours l'apanage des confesseurs de la foi. M. du Coin était le plus soupçonné, aussi fut-il soumis à de plus rudes épreuves ; mais l'esprit de foi dont il se couvrait comme d'un bouclier le rendait invulnérable. Sous cette égide tutélaire, il repoussait avec facilité tous les traits qu'on lançait contre lui. On voulait absolument parvenir à obtenir de lui quelques éclaircissements sur les moyens de correspondance établis pour communiquer avec le Saint-Père ; on multipliait à ce sujet les interrogatoires ; on s'efforçait de l'embarrasser et de lui arracher quelque aveu en l'accablant des questions les plus astucieuses. Lorsqu'on s'aperçut qu'il fallait renoncer à l'ébranler ainsi, on voulut essayer de l'effrayer par des menaces. On lui annonce donc qu'on va le charger de chaînes : Vous me ferez plaisir, répond avec calme ce héros chrétien, car alors elles seront pour moi la livrée de J. C., et je considérerai comme un très grand bonheur de les porter pour la gloire de son nom. On chercha à lasser sa patience ; on le tint assez longtemps claquemuré dans un cachot
% Source : PDF original, vue 125, côté droit, page manuscrite 251.
de six pieds en carré. Alors notre saint ami prenant pour modèles ces saints des temps antiques qui dans les siècles de ferveur peuplèrent les déserts de la Thébaïde, s'imposa lui-même une règle à laquelle il se montra constamment fidèle. L'heure du lever, du coucher, de la méditation, de la prière vocale, des repas, de la promenade même malgré les bornes étroites dans lesquelles il se trouvait renfermé, tout était prévu et fut exécuté avec la plus scrupuleuse exactitude. Ses persécuteurs étaient vaincus ; il ne leur était pas donné de concevoir quelle force une âme entièrement soumise et résignée à la volonté du souverain maître peut recevoir d'en haut. Ils eurent recours à un dernier moyen pour abattre ce grand courage : ces hommes tout charnels, ne connaissant d'autre mobile de conduite que les passions qui les dominaient, se persuadent en conséquence que le sentiment qui fait agir ce fervent serviteur de Dieu est un sentiment d'orgueil ; remplis de cette pensée ils prennent la résolution de l'humilier, ce sera le moyen de triompher de cette résistance jusqu'alors invincible. À cette fin ils le font changer de quartier, ils placent cet homme dont l'âme est si pure et si loyale au milieu de cette classe hideuse et dégoûtante que le crime a rassemblée dans les prisons ; mais le juste vit de la foi et se rit de leurs misérables efforts. Comme un rocher que les flots irrités frappent sans relâche à coups redoublés de leurs vagues furieuses sans parvenir à l'ébranler, du Coin toujours impassible, tranquille avec sa conscience, saura faire tourner à la gloire de son divin maître cette rigueur mesquine dont les auteurs auraient dû rougir ; il rassemble ces malheureux auxquels on a voulu l'associer. Le voilà au milieu d'eux, un catéchisme à la main ; il leur explique la loi de Dieu, ils écoutent avec attention et respect ; l'ascendant que lui donne sa
% Source : PDF original, vue 126, côté gauche, page manuscrite 252.
haute vertu porte la conviction dans tous les esprits, et la bonté avec laquelle il se met à la portée de tous ébranle et touche tous les cœurs. Ces hommes le rebut de la société rentrent en eux-mêmes et s'empressent de revenir à la pratique de ces devoirs dont l'oubli est la première cause de leurs malheurs.

Tout dans ce monde trouve enfin un terme, même les rigueurs de la persécution. Après avoir épuisé toutes les ressources de leur art séducteur, ces hommes de la police se lassèrent de travailler en vain ; ils cessèrent d'interroger nos amis, ils se relâchèrent un peu de ce régime sévère auquel ils les avaient d'abord condamnés. Nos prisonniers purent enfin recevoir des nouvelles de leurs familles, mais le gouvernement de Napoléon ne rendait guère ses victimes à la liberté ; nos captifs passèrent donc plus de trois années sous les verroux ; Dieu leur réserva pour la dernière année une grande consolation, ils purent se voir, et furent réunis dans la même prison. Là comme ailleurs, et peut-être mieux qu'ailleurs on comprend les douceurs qui sont attachées à l'\textit{Ecce quam bonum et quam jucundum}\footnote{(Note du transcripteur) Littéralement : « Voici comme il est bon et comme il est agréable ». Psaume 132 (133), verset 1.}.

C'était sans doute un adoucissement à leurs peines, mais quant à la liberté le souverain maître qui voulait que la Congrégation de Lyon fût constamment unie avec le St Siège dans la persécution comme dans la prospérité, ne permit pas qu'ils en fussent mis en possession avant le chef de l'Église avec lequel et pour lequel ils étaient persécutés. Ce ne fut donc que lorsque les puissances du Nord envoyées en France par le Dieu des armées pour y accomplir ses desseins se furent emparées de Paris, que nos amis furent enfin délivrés.

Avant de nous occuper de cette époque intéressante, il nous reste
% Source : PDF original, vue 126, côté droit, page manuscrite 253, avant le chapitre onzième.
à suivre la marche de notre Congrégation pendant ces tristes années de persécution ; c'est ce qui fera l'objet du chapitre suivant.
